%%
%% Author: shoupu_wan
%% 7/19/18
%%
\section{Queue}\label{sec:queue}

A queue is a stateful collection that is meant to be modified with push and pop operations
frequently.

Using a queue can free us from keeping track of the iteration progress explicitly and thus we can
iterate the same data structure in various functions conveniently as illustrated in the example
below.
\lstinputlisting[
caption={Illustration for using queue for stateful iteration across multiple functions}]
{data-structure/JsonTokenizer_atomic.py}

\section{Stack}\label{sec:stack}
A stack is a mimick of the stack used by operating system to manage stack memory.

\section{Stack \textit{vs.} Queue}
\begin{itemize}
    \item DFS \textit{vs.} BFS \label{Evernote:DFS vs BFS}
    \item Graph exploring patterns
\end{itemize}


\label{Evernote:Versioned Stack algorithm}
You need to implement a versioned stack, i.e. version of stack will increment after each push/pop. Apart from push/pop, implement a method print(n) which would print stack state corresponding to version 'n'. For example:
-> initially stack is empty.
-> Version 1: 11 is pushed
-> Version 2: 8 is pushed
-> version 3: pop. only 11 left
-> Version 4: 15 is pushed
....
And so on.

Print(n) should print state at version 'n'.

Here 1 should print 11, 2 should print 8, 11...
All methods should be as efficient as possible.
